% The LoCoBench generation pipeline: five prompted stages, then three validators.
%
% Extracted from submit.tex so the figure has a single source and can be exported
% standalone by scripts/make_bench_pipeline_png.sh. The validators are V1, V3 and
% V4 by design -- those are the three that gate admission, and the numbering
% matches locobench/validate.py's threshold keys (v1_coupling, v1_sense,
% v3_min_score, v4_coverage). There is deliberately no V2 box.
%
% The committee that runs the validators EXCLUDES the generating model, which is
% why the two left-hand annotations name different model counts.
% ---------------------------------------------------------------------------
\begin{figure}[tbp]
\centering
\resizebox{\textwidth}{!}{%
\begin{tikzpicture}[font=\small,
  st/.style={draw,rounded corners=1.5pt,align=center,inner sep=3.5pt,minimum height=9mm,fill=blue!5},
  vd/.style={draw,rounded corners=1.5pt,align=center,inner sep=3.5pt,minimum height=9mm,fill=orange!10},
  fin/.style={draw,rounded corners=1.5pt,align=center,inner sep=3.5pt,minimum height=9mm,fill=green!12},
  ar/.style={-{Stealth[length=1.9mm]},thick}, node distance=3.4mm]
  \node[st] (t) {topic};
  \node[st,right=of t] (p1) {\textbf{P1}\\question};
  \node[st,right=of p1] (p2) {\textbf{P2}\\$26$ claims};
  \node[st,right=of p2] (p3) {\textbf{P3}\\relation plan};
  \node[st,right=of p3] (p4) {\textbf{P4}\\base response};
  \node[st,right=of p4] (p5) {\textbf{P5}\\$\times 4$ rungs};
  \foreach \a/\b in {t/p1,p1/p2,p2/p3,p3/p4,p4/p5} \draw[ar] (\a) -- (\b);
  \node[vd,below=9mm of p3] (v1) {\textbf{V1} recover\\relations};
  \node[vd,right=of v1] (v3) {\textbf{V3} prose\\audit};
  \node[vd,right=of v3] (v4) {\textbf{V4} claim\\coverage};
  \node[fin,right=of v4] (adm) {admit family\\(all $5$ rungs)};
  \draw[ar] (p5.south) -- ++(0,-3mm) -| (v1.north);
  \draw[ar] (v1) -- (v3); \draw[ar] (v3) -- (v4); \draw[ar] (v4) -- (adm);
  \node[left=1mm of v1,font=\scriptsize\itshape,align=right,text width=15mm] {committee\\($3$ models)};
  \node[left=1mm of t,font=\scriptsize\itshape,align=right,text width=15mm] {generator\\($1$ model)};
\end{tikzpicture}}
% scripts/make_bench_pipeline_png.sh defines \bpnocaption to render this figure
% for standalone export, where surrounding prose supplies the explanation.
% Undefined in the paper, so the caption below is typeset normally there.
\ifdefined\bpnocaption\else
\caption{The generation pipeline. Five stages produce a family of five rungs from one topic; three
  validators, run by a committee that \emph{excludes} the generating model, gate admission. A family is
  admitted whole or rejected whole (Def.~\ref{def-ladder}).}
\label{fig:bench-pipeline}
\fi
\end{figure}
